\section{Resumen}

Esta tesis estudia el seguimiento de precios en un creador automático de mercado (AMM) sintético de producto constante y sin comisiones, frente a una referencia exógena de un mercado centralizado (CEX). Separa las consecuencias algebraicas de la regla de intercambio de los efectos que dependen del comportamiento de los agentes. En una piscina aislada, un arbitrajista limitado a una fracción c de la reserva de entrada puede reducir el error logarítmico fuera de la banda hasta en dos veces el logaritmo natural de uno más c por tick. Para un choque logarítmico de 0.20 y una banda de 0.003, el tiempo idealizado hasta la banda es de 20 ticks con c = 0.005, cuatro con c = 0.03 y dos con c = 0.10. Este cálculo explica la dirección del resultado de arbitraje configurado; la simulación mide desviaciones bajo ruido y otras reglas.

En 13,500 ejecuciones con 300 semillas Monte Carlo, el escenario de capacidad alta tiene un área media de exceso de error de 0.10, frente a 2.36 sin arbitraje. El resultado depende de la regla y la grilla de parámetros especificadas, y el signo del efecto está incorporado mecánicamente en esa regla. El disparador actual de retiro de liquidez se activa en todas las corridas de la condición adaptativa, normalmente en el primer tick. Su aparente aumento de la respuesta máxima ocurre solo con la menor profundidad y órdenes absolutas fijas; desaparece al escalar las órdenes con las reservas. La métrica original de desplazamiento máximo no mide el sobrepaso respecto de la nueva referencia, por lo que esta sensibilidad no demuestra una amplificación general de liquidez.

El hallazgo del simulador con mayor potencial de transferencia se refiere a la persistencia: el 31.2 por ciento de las trayectorias que volvió a la banda de no arbitraje salió de ella después; con un factor de profundidad de 0.5, la proporción fue 98.8 por ciento. Al definir medidas de monitoreo conviene distinguir el primer retorno del seguimiento sostenido. El mapa reducido de dos estados mejora la persistencia en algunas condiciones, pero no supera a un AR(1) escalar. Un mapa sensible a la banda reduce cerca de 1.7 por ciento el error recursivo de reproducción, concentrado en una condición y usando la trayectoria futura realizada del CEX; no es un resultado de pronóstico en tiempo real. La contribución es un estudio de simulación reproducible y centrado en mecanismos que distingue resultados algebraicos de sensibilidades dependientes de reglas, y muestra dónde fallan sus modelos reducidos. La reproducción de eventos con datos de mercado sincronizados y costos realistas queda para trabajo futuro.

\textbf{Palabras clave:} creador automático de mercado; arbitraje; seguimiento de precios; retiro de liquidez; simulación basada en agentes; experimento Monte Carlo
